\section{Related work}
\label{sec:related}

\paragraph{Semantic caching for language models.}
Systems that reuse model outputs by embedding similarity have moved quickly from prototype
to production. \mbox{GPTCache} established the pattern of encoding a query, retrieving its
nearest stored neighbor, and serving the neighbor's response above a similarity
threshold~\cite{bang2023gptcache}. Later work improved the hit rate and the threshold
itself: \mbox{MeanCache} personalizes the decision with federated
learning~\cite{gill2024meancache}, category-aware caching sets different thresholds for
different query types~\cite{wang2025categoryaware}, and domain-specific embeddings sharpen
the similarity signal with synthetic fine-tuning data~\cite{gill2025advancing}. The work
closest to ours is \mbox{vCache}, which learns a per-prompt threshold online from
correctness feedback and bounds the error rate with a statistical
guarantee~\cite{schroeder2025vcache}. These systems share an objective: maximize reuse
subject to keeping answers correct, where correctness is judged on the served queries.

\paragraph{What has not been measured.}
The literature reports hit rates, latency, and, in the case of verified caching, a bound on
error~\cite{schroeder2025vcache,gill2024meancache}. It does not chart how the benign
false-hit rate moves across embedding models and domains, it does not follow a false hit
into the quality of the downstream task, and it does not turn the cost of a false hit into
an operating point. Our contribution is orthogonal to building a better cache: we measure
the failure that every similarity cache shares and calibrate against its measured cost.
The one study that probes these caches adversarially shows that an attacker can craft
collisions that force false hits~\cite{zhang2026keycollision}; we ask the complementary
question of how often the same failure occurs without an adversary.

\paragraph{Embeddings and their evaluation.}
The encoders we audit are standard sentence embedding models trained with contrastive
objectives: Sentence-BERT and its \mbox{MPNet} backbone~\cite{reimers2019sentencebert,song2020mpnet},
the weakly supervised \mbox{E5} models~\cite{wang2022e5}, and the \mbox{BGE} models from the
\mbox{C-Pack} release~\cite{xiao2023cpack}, all of which rank highly on the Massive Text
Embedding Benchmark~\cite{muennighoff2023mteb}. Our labeled domains come from human-annotated
paraphrase and duplicate-question corpora~\cite{wang2019glue,dolan2005mrpc,zhang2019paws}
rather than from semantic-similarity regression~\cite{cer2017sts}, because the question a
cache must answer is binary equivalence, not graded similarity.

\paragraph{Measurement methodology.}
We follow the empirical-rigor practices that the evaluation literature has converged on:
confidence intervals from the bias-corrected and accelerated
bootstrap~\cite{efron1987bootstrap,efron1993bootstrap}, multiple-comparison correction by
controlling the false discovery rate~\cite{benjamini1995fdr}, and significance testing
chosen for the structure of the data rather than by default~\cite{dror2018hitchhiker}. We
document the data and the models with a datasheet and model cards in the
spirit of~\cite{gebru2021datasheets,mitchell2019modelcards}.
